% Generated draft: requires mathematical and source review.
\documentclass[11pt]{article}
\usepackage[margin=25mm]{geometry}
\usepackage{fontspec}
\setmainfont{Latin Modern Roman}
\usepackage{amsmath,amssymb,mathtools,unicode-math}
\setmathfont{Latin Modern Math}
\usepackage{enumitem,needspace,xurl,hyperref}
\setlength{\parindent}{0pt}
\setlength{\parskip}{5pt}
\setlength{\emergencystretch}{3em}
\newcommand{\N}{\mathbb{N}}
\newcommand{\Z}{\mathbb{Z}}
\newcommand{\Q}{\mathbb{Q}}
\newcommand{\R}{\mathbb{R}}
\newcommand{\C}{\mathbb{C}}
\newcommand{\F}{\mathbb{F}}
\newcommand{\A}{\mathbb{A}}
\newcommand{\PP}{\mathbb{P}}
\newcommand{\E}{\mathbb{E}}
\newcommand{\eps}{\varepsilon}
\newcommand{\dd}{\,\mathrm{d}}
\begin{document}
\section*{UnsolvedMath ID: 701676}
\textbf{Source problem code:} catalogue entry
\subsection{Definitions and mathematical statement}
For a prime \(p\) and an integer \(k\geq 2\), let
\[
R_{p^k}=\{G: G\models \forall x\,x^{p^k}=1\text{ and }\forall x,y,z\,[[x,y],z]=1\},
\]
where \([x,y]=x^{-1}y^{-1}xy\). Thus \(R_{p^k}\) is the variety of groups of nilpotency class at most \(2\) and exponent dividing \(p^k\). These bounds are the conventional interpretation of the source's phrase ``variety of class \(2\) nilpotent groups of exponent \(p^k\).'' Write \(G' = \langle [x,y]:x,y\in G\rangle\) for the derived subgroup. For a finite group \(H\), \(\exp(H)\) is the least positive integer \(e\) such that \(h^e=1\) for every \(h\in H\).

A group quasi-identity is a universally quantified implication
\[
\forall x_1,\ldots,x_n\;\left(\bigwedge_{j=1}^{m}u_j=v_j\ \Longrightarrow\ u=v\right),
\]
where all displayed expressions are group terms and the conjunction may be empty. A quasivariety is a class of groups axiomatized by quasi-identities. Let \(\mathsf Q(F)\) denote the smallest quasivariety containing a group \(F\), equivalently the class of groups satisfying every quasi-identity valid in \(F\).

A basis of quasi-identities for a quasivariety \(\mathcal K\) is a set \(B\) whose group models are exactly \(\mathcal K\). It is independent if, for every \(\varphi\in B\), there exists a group satisfying \(B\setminus\{\varphi\}\) and failing \(\varphi\). This defines independence over the group axioms. A relative basis within \(R_{p^k}\) instead has \(\mathcal K=\{G\in R_{p^k}:G\models B\}\), and relative independence requires each witnessing group to belong to \(R_{p^k}\). The supplied record does not specify which meaning of basis and independence is intended; that qualification remains unresolved.

\textbf{Question (A. I. Budkin, 2022).} Is it true that, for every prime \(p\) and every integer \(k\geq 2\), there exist infinitely many pairwise distinct subquasivarieties \(\mathcal K_i\subseteq R_{p^k}\), indexed by \(i\geq 1\), and finite groups \(F_i\in R_{p^k}\) such that
\[
\mathcal K_i=\mathsf Q(F_i),\qquad \exp(F_i')=p^k,
\]
and each \(\mathcal K_i\) has no independent basis of quasi-identities? The basis convention must be resolved as described above. The equality \(\exp(F_i')=p^k\) retains the source's stated exponent requirement as an exact exponent. The universal quantifier includes \(p=2\); the source supplies no exception for that prime.
\subsection{Short English statement}
For every prime \(p\) and integer \(k\geq 2\), Budkin asks whether there are infinitely many distinct quasivarieties within \(R_{p^k}\), each generated by a finite group whose derived subgroup has exponent exactly \(p^k\), and each lacking an independent basis of quasi-identities. The source leaves unresolved whether the basis and independence are meant over all groups or relative to \(R_{p^k}\).
\subsection{Sources}
\begin{itemize}
\item Local UnsolvedMath registry, permanent ID \(701676\), exact source problem code \texttt{catalogue entry}, title \emph{Kourovka Notebook Problem 20.18}; supplied upstream revision \texttt{d371ee917ac9647fbb06cd6fc529deec6f3dda82}. The registry labels the problem open; this is a supplied status, not an independently verified mathematical assessment.
\item \emph{The Kourovka Notebook (Unsolved Problems in Group Theory)}, Problem 20.18, proposed by A. I. Budkin in 2022. The supplied background identifies the 20th issue (2022) and the 21st issue edition (2026), \texttt{arXiv:1401.0300v48}. The supplied collection description names E. I. Khukhro and V. D. Mazurov as editors of Issues 1--21 (1965--2026). These bibliographic details are reproduced from the record without independent source consultation.
\end{itemize}
\end{document}
